% Rebuttal to Reviewer 9v2E
% ACM-BCB 2026 Submission: Structure-Aware Prediction of PROTAC-Mediated
% Protein Degradability via Graph Neural Networks
%
% Compile: pdflatex rebuttal_9v2E.tex && pdflatex rebuttal_9v2E.tex

\documentclass[11pt]{article}
\usepackage[margin=1in]{geometry}
\usepackage{xcolor}
\usepackage{booktabs}
\usepackage{enumitem}
\usepackage{hyperref}
\usepackage{amsmath}

% Styling
\definecolor{revblue}{RGB}{0,70,170}
\newcommand{\reviewerbox}[2]{%
  \noindent\fcolorbox{black}{gray!10}{\parbox{\dimexpr\linewidth-2\fboxsep-2\fboxrule}{%
    \textbf{#1}: #2}}%
  \vspace{4pt}%
}
\newcommand{\resp}{\noindent\textbf{Response:}~}
\newcommand{\change}[1]{\textcolor{revblue}{[Paper reference: #1]}}

\setlength{\parskip}{6pt}
\setlength{\parindent}{0pt}

\begin{document}

\begin{center}
{\Large\bfseries Response to Reviewer 9v2E}\\[6pt]
{\large Structure-Aware Prediction of PROTAC-Mediated Protein Degradability\\via Graph Neural Networks}\\[6pt]
ACM-BCB 2026 Submission\\[4pt]
\textit{Rating: 4 --- Borderline Accept}
\end{center}

\bigskip

We thank Reviewer 9v2E for the constructive feedback. We have substantially revised the manuscript; all changes are \textcolor{revblue}{highlighted in blue} in the revised paper. Line numbers below refer to the revised \texttt{paper.tex}.

\bigskip
\hrule
\bigskip

%% ====================================================================
%%  C1
%% ====================================================================

\reviewerbox{C1}{The dataset is relatively small (3,101 samples, 155 unique proteins, 10 E3 ligases). The target-unseen effective diversity is much closer to the number of unique proteins, contributing to seed sensitivity.}

\resp We fully agree. We have added a prominent ``Dataset characteristics and limitations'' paragraph \emph{before} any results are presented:

\begin{quote}\itshape
``PROTAC-8K contains 3,101 labeled samples spanning only 155 unique protein targets\ldots\ The E3 ligase distribution is highly imbalanced: CRBN accounts for 62\% of samples, VHL for 35\%, and the remaining eight E3s combined represent only 3\%.\ These constraints fundamentally bound generalization performance\ldots''
\end{quote}

We also note that PROTAC-8K is the only publicly available benchmark of this scale, and we have no external validation dataset---both stated as explicit limitations.

\change{Lines 431--436 (new ``Dataset characteristics and limitations'' paragraph, \S4 Results, before any performance numbers); Limitation item~1 at line 766 (small dataset); Limitation item~2 at line 767 (E3 imbalance); Limitation item~6 at line 771 (no external validation).}

\bigskip

%% ====================================================================
%%  C2
%% ====================================================================

\reviewerbox{C2}{Structural validation is limited. Predictions are not supported by docking, ternary-complex modeling, MD, or free energy calculations.}

\resp We agree this is a limitation. It is now discussed explicitly in the Limitations paragraph:

\begin{quote}\itshape
``Predictions rely on static AlphaFold structures without physics-based validation (docking, MD, free energy calculations). The BRD4 case study (K374) provides one anecdotal positive control.''
\end{quote}

We note an inherent tension: ternary-complex modeling requires knowing the PROTAC linker geometry, which is unavailable in our pre-synthesis setting. Post-hoc structural validation against experimentally determined ternary complexes (where available) is an important future direction.

\change{Limitation item~5 at line 770 (``No structural validation''); BRD4 case study (\S4.5, lines 686--703) provides partial structural validation via lysine attention weight analysis on a known target.}

\bigskip

%% ====================================================================
%%  C3
%% ====================================================================

\reviewerbox{C3}{No compelling biological case study showing why a specific protein is predicted degradable, which lysines drive the prediction, or how predictions align with literature.}

\resp We have added a new subsection ``Case Study: BRD4 Degradability'' (\S4.5) that provides exactly this analysis:

\begin{itemize}[topsep=2pt,itemsep=1pt]
  \item BRD4 (O60885) is predicted highly degradable with CRBN (score: 0.87). \change{Line 690.}
  \item Lysine attention weights identify \textbf{K374} (attention: 0.12, SASA: 94\,\AA$^2$) as the top residue---located in a disordered, solvent-exposed linker between bromodomains. \change{Lines 691--697.}
  \item \textbf{Literature validation}: mass spectrometry studies by Zengerle et al.\ (2015) confirm K374 as a major ubiquitination site following JQ1-based PROTAC treatment. \change{Lines 699--702.}
  \item We acknowledge this is a single successful case and that systematic validation across more targets would strengthen confidence. \change{Line 703.}
\end{itemize}

\change{New \S4.5 ``Case Study: BRD4 Degradability'' at lines 686--703.}

\bigskip

%% ====================================================================
%%  Additional concerns (implicit in review)
%% ====================================================================

\reviewerbox{Additional}{The key target-unseen result is unstable; the paper emphasizes the best seed over the multi-seed result; some biological and E3 generalization claims are only partly supported.}

\resp All three points are now addressed:

\begin{enumerate}[topsep=2pt,itemsep=1pt]
  \item \textbf{Instability addressed with 6-seed validation}: We trained 3 additional seeds (789, 1011, 1213), yielding a 6-seed mean of 0.603$\pm$0.097---marginally below the gradient boosting baseline (0.607, bootstrap $p=0.556$). We explicitly recommend ensembling across $\geq$6 seeds for deployment.
  \change{Lines 619--625 (``Six-seed extended validation'' paragraph); Figure~4 at line 632; Table~1 at lines 441--458 (both 3-seed and 6-seed rows).}

  \item \textbf{Best seed de-emphasized}: The multi-seed average is now the primary result throughout; best seed (0.7449) is consistently marked as secondary and parenthetical.
  \change{Abstract line 65 (``0.646$\pm$0.124 AUROC\ldots (best seed: 0.7449)''); Introduction line 135 (``3-seed mean of 0.646\ldots and 6-seed mean of 0.603''); Table~1 line 444 (``Best Seed'' as separate column); Results line 460--462 (``multi-seed averages as primary results\ldots best-seed performance should not be interpreted as typical expectation''); Conclusion line 783 (leads with average).}

  \item \textbf{E3 claims narrowed}: Changed throughout from ``E3 generalization'' to ``CRBN$\to$VHL transfer.'' VHL failure (0.396 AUROC) and root cause analysis now prominently discussed. Rare E3 generalization explicitly stated as unsupported.
  \change{Abstract line 65 (``CRBN$\to$VHL E3-unseen transfer''); Section heading line 635 (``E3 generalization is limited to major ligases''); Root cause analysis lines 640--648; Per-E3 Table~5 at line 654; Per-E3 Figure~5 at line 668; Conclusion line 783 (``CRBN$\to$VHL transfer (0.811 AUROC)'').}
\end{enumerate}

\bigskip
\hrule
\bigskip

\section*{Summary of Changes Relevant to Reviewer 9v2E}

\begin{center}
\small
\begin{tabular}{@{}p{3.5cm}p{8cm}@{}}
\toprule
\textbf{Concern} & \textbf{Revision (lines in \texttt{paper.tex})} \\
\midrule
Small dataset (C1) & \S4 L431--436 (dataset limitations paragraph); Limitations L766, L767, L771 \\
No structural validation (C2) & Limitation~5 L770; \S4.5 L686--703 (BRD4 partial validation) \\
No case study (C3) & New \S4.5 L686--703 (BRD4 with K374 literature validation) \\
Seed instability & \S4.3 L619--625 (6-seed validation); Table~1 L441--458; Figure~4 L632 \\
Best seed emphasis & Abstract L65; Introduction L135; Table~1 L444; \S4.2 L460--462; Conclusion L783 \\
E3 claims too broad & Abstract L65; \S4.3 L635--648; Table~5 L654; Figure~5 L668; Conclusion L783 \\
\bottomrule
\end{tabular}
\end{center}

\end{document}
